From now on, we let $J_k$ be a job of $\tau_k$ that has the worse-case response time. As in~[6], we extend the beginning of the level-$k$ busy period from $r_k$ (the release time of $J_k$), to an earlier time instant $t_0$, which is defined as the earliest time instant before $r_k$, such that at any time instant $t \in [t_0, r_k)$ all processors are occupied by tasks with higher priority than $\tau_k$. If there is no such a time instant, we set $t_0 = r_k$. Using this, the \textit{level-k busy period} is defined as the time interval $[t_0, f_k)$, and we define $\varphi = r_k - t_0$, which is the time span by which the busy period extends to the left, as shown in Figure~2.

This definition of $t_0$ is chosen to impose a bound on the number of tasks contributing carry-in, since at time point $t_0 - 1$, there have to be strictly less than $M$ higher priority tasks active. We state that property as the following lemma:

\textbf{Lemma 1.} \ \textit{There are at most $M - 1$ tasks having carry-in, and for each task $\tau_i$, the carry-in is at most $C_i - 1$.}

\textit{Proof:} By discussion above, similar to~[6]. \hfill $\blacksquare$